% DOCUMENT CLASS %
\documentclass[a4paper, 12pt]{extarticle}   % for A4 Paper, 12pt fontsize and article spacing
\usepackage[left=1.5cm,right=1.5cm,top=2cm,bottom=2cm]{geometry}    % Make the page a little bigger

% PACKAGES%
\usepackage{amsfonts}   % For a basic mathfont (many will be replaced by stix)
\usepackage{amsmath}    % For basic math symbols
\usepackage{bbm}        % For \mathbbm (better version of \mathbb)
\usepackage{mathrsfs}   % For \mathscr
\usepackage{hyperref}   % For footnotes
\usepackage{csquotes}   % For \textquote{}
\usepackage{IEEEtrantools}  % For better alignments
\usepackage{tikz}   % For a macro
\usetikzlibrary{arrows.meta, positioning, calc}	% For protocol diagrams
\usepackage{listings} % For code
\usepackage{xcolor}   % For ... code
\usepackage{microtype}  % For better paragraph formatting
\usepackage{amsthm}   % For custom environments

% FONT CONFIGURATION %
\usepackage[T1]{fontenc}
\usepackage{palatino}

% BASIC SETS %
\newcommand*{\R}{\mathbbm R}    % Set of real numbers
\newcommand*{\N}{\mathbbm N}    % Set of natural numbers (beginning at 0)
\newcommand*{\Z}{\mathbbm Z}    % Set of integers
\newcommand*{\Q}{\mathbbm Q}    % Set of rational numbers

% MACROS %
\newcommand*{\qedbox}{\null\nobreak\hfill\ensuremath{\square}} % Macro for a QED-Box
% for a funny ¯\_(ツ)_/¯-Emoji
\newcommand{\shrug}[1][]{%
    \begin{tikzpicture}[baseline,x=0.8\ht\strutbox,y=0.8\ht\strutbox,line width=0.125ex,#1]
    \def\arm{(-2.5,0.95) to (-2,0.95) (-1.9,1) to (-1.5,0) (-1.35,0) to (-0.8,0)};
    \draw \arm;
    \draw[xscale=-1] \arm;
    \def\headpart{(0.6,0) arc[start angle=-40, end angle=40,x radius=0.6,y radius=0.8]};
    \draw \headpart;
    \draw[xscale=-1] \headpart;
    \def\eye{(-0.075,0.15) .. controls (0.02,0) .. (0.075,-0.15)};
    \draw[shift={(-0.3,0.8)}] \eye;
    \draw[shift={(0,0.85)}] \eye;
    % draw mouth
    \draw (-0.1,0.2) to [out=15,in=-100] (0.4,0.95); 
    \end{tikzpicture}
}

% BASIC CONFIGURATION %
\pagestyle{plain}
\allowdisplaybreaks

% CODE CONFIGURATIONS %
\definecolor{codepurple}{HTML}{7f0055}
\definecolor{comment}{RGB}{0,128,0}
\definecolor{number}{RGB}{9,136,90}
\definecolor{string}{RGB}{163,21,21}

\lstset{
    language=python,
    emph={*, do, od, for, if, fi, then, else, to, def, in, forall, exists},
    firstnumber=1,
    basicstyle=\small\ttfamily,
    breaklines=true,
    breakatwhitespace=false,
    frame=tb,
    captionpos=b,
    %float=t,
    %floatplacement=t,
    abovecaptionskip=0.5cm,
    numbers=left,
    numberstyle=\tiny,
    keywordstyle=\bfseries\color{codepurple},
    emphstyle=\bfseries\color{codepurple},
    mathescape=true,
    keepspaces=true,
    showspaces=false,
    showstringspaces=false,
    showtabs=false,
    stringstyle=\color{string},
    commentstyle=\itshape\color{comment},
    upquote=true,
    texcl=true
}

% EXERCISE ENVIRONMENT %
\theoremstyle{definition}
\newtheorem{exercise}{Exercise}[section]
\setcounter{section}{1} % Don't use \section and this works naturally

% TITLE CONFIGURATION %
\newcommand{\ex}{0}

\title{Data Visualization \\ \small Exercise Sheet \ex}
\date{Winter Term 2025/26}
\author{
    Henri Heyden \\%
    \small stu240825%
    \and
    Julia Schröder \\%
    \small stu252018%
    \and
    Rebecca Ziebell \\%
    \small stu250567%
}

\usepackage{multirow}
\renewcommand{\ex}{6}
\setcounter{section}{\ex}
\newcommand{\Prop}{\text{Pr}}

\begin{document}
    \maketitle

    \begin{exercise}\hfill
        \begin{center}
        \scalebox{0.4}{
        \begin{tabular}{|l|l|l|l|l|l|l|}
        \hline
        \multicolumn{7}{|c|}{Tabular Data Visualization Idioms} \\
        \hline
        Name & Data & Derived Data & Marks & Channels & Tasks & Scalability \\
        \hline
        Scatterplot &
        2 quantitative values &
        -- &
        Points &
        x/y position, color, size &
        Correlation, clusters, outliers &
        Hundreds of items \\
        \hline
        Bar Chart &
        1 categorical key, 1 quantitative value &
        -- &
        Bars (lines) &
        Length, aligned position &
        Compare values, lookup &
        Dozens--hundreds of categories \\
        \hline
        Stacked Bar Chart &
        2 categorical keys, 1 quantitative value &
        -- &
        Stacked bars &
        Length, color hue &
        Part-to-whole comparison &
        Main key: 100s, stacked key: $\sim$10 \\
        \hline
        Normalized Stacked Bar Chart &
        2 categorical keys, 1 quantitative value &
        Normalization to 100\% &
        Stacked bars &
        Relative length, color &
        Relative part-to-whole &
        Similar to stacked bars \\
        \hline
        Line / Dot Chart &
        1 ordered key, 1 quantitative value &
        -- &
        Points + lines &
        Aligned position &
        Trend detection &
        Hundreds of items \\
        \hline
        Indexed Line Chart &
        1 ordered key, 1 quantitative value &
        Indexed values &
        Points + lines &
        Relative position &
        Change over time &
        Same as line chart \\
        \hline
        Streamgraph &
        1 categorical key, 1 ordered key, 1 value &
        Layer geometry &
        Filled areas &
        Height, color &
        Evolution over time &
        100s time steps, 100s categories \\
        \hline
        Gantt Chart &
        1 categorical key, 2 quantitative values &
        Duration &
        Line segments &
        Horizontal position, length &
        Temporal overlaps &
        Dozens of items \\
        \hline
        Slopegraph &
        2 quantitative values &
        Change magnitude &
        Points + lines &
        Vertical position, slope &
        Change and rank comparison &
        Dozens of items \\
        \hline
        Heatmap &
        2 categorical keys, 1 quantitative value &
        -- &
        Cells &
        Color &
        Patterns, clusters &
        Up to millions of cells \\
        \hline
        Clustered Heatmap &
        Heatmap data &
        Clustering hierarchies &
        Cells + tree edges &
        Color, position &
        Assess clustering quality &
        Limited by clustering \\
        \hline
        Pie Chart &
        1 categorical key, 1 quantitative value &
        -- &
        Sectors &
        Angle / arc length &
        Part-to-whole (few categories) &
        Few categories only \\
        \hline
        SPLOM &
        Multiple quantitative attributes &
        -- &
        Points &
        x/y position &
        Pairwise correlations &
        $\sim$10 attributes \\
        \hline
        Parallel Coordinates &
        Multiple quantitative attributes &
        -- &
        Polylines &
        Position along axes &
        Multivariate patterns &
        Dozens of attributes \\
        \hline
        \end{tabular}
        }
    \end{center}
    \end{exercise}

    \begin{exercise}
    \end{exercise}

    \begin{exercise} \hfill
        \begin{itemize}
            \item[a)]
                \begin{itemize}
                    \item Pros \\
                    Connections between the data are easier to understand because of the arrows. \\
                    Different colors for different aspects. \\
                    
                    \item Cons \\
                    Too many arrows result in the diagramm being confusing. \\
                    It is difficult to understand the meanings behind those many arrows. \\
                    Without a good overview, it is difficult to compare the data. \\
                \end{itemize}
            \item[c)] See \autoref{fig:3a}, \autoref{fig:3b}, \autoref{fig:3c}, and \autoref{fig:3d}.
            \item[d)]
                Pie Charts
                \begin{itemize}

                    \item Pros \\
                    Intuitive for part-to-whole judgements within a single year.
                    Easy to see which categories dominate in 2016 or 2017 individually.
                    Familiar to most audiences. 

                    \item Cons \\
                    Very poor for comparison across years: viewers must mentally compare angles/areas between two seperate pies.
                    Small differences are hard to detect.
                    Ordering of slices is often unclear or inconsisten.
                    Does not emphasize change, only composition. 

                    \item What is highlighted/not highlighted \\
                    Highlights relative proportions within one year.
                    Does not highlight change between 2016 and 2017.
                    
                \end{itemize}
                
                Stacked Bar Chart
                \begin{itemize}

                    \item Pros \\
                    Still supports part-to-whole comparison per year.
                    Easier to compare totals and large segments than pie charts.
                    Uses aligned axes, improving quantitative perception. 

                    \item Cons \\
                    Only the bottom segment is well aligned, upper segments are hard to compare across years.
                    Change in individual categories is not clearly emphasized unless  ordering is carefully chosen.
                    Can become cluttered with many categories. 

                    \item What is highlighted/not highlighted \\
                    Highlights overall composition and total size.
                    individual category changes (except baseling segment) are hard to compare.
                \end{itemize}

                Bar Chart
                \begin{itemize}

                    \item Pros \\
                    Very good for direct comparison between 2016 and 2017.
                    Aligned bar lengths make differences easy to perceive.
                    Supports accurate value lookup.
                    Ordering categories by magnitude or change improves readability.

                    \item Cons \\
                    Loses the explicit part-to-whole perspective.
                    Requires more horizontal space.
                    Can become visually dense with many categories. 

                    \item What is highlighted/not highlighted \\
                    Highlights absolute differences between years per category.
                    Does not emphasize overall composition. 

                \end{itemize}

                Slopegraph
                \begin{itemize}

                    \item Pros \\
                    Specifically designed to show change between two time points.
                    Makes direction and magnitude of change immediately visible.
                    Very effective for ranking changes and identifying increases/decreases.
                    Compact and expressive. 

                    \item Cons \\
                    Exact values are harder to read.
                    Can become cluttered if there are many categories.
                    Requires careful labeling to avoid ambiguity. 

                    \item What is highlighted/not highlighted \\
                    Strongly highlights change and rank shifts.
                    Less suitable for precise value reading or part-to-whole analysis. 
            
                \end{itemize}
        \end{itemize}
        
    \end{exercise}

\end{document}